\documentclass[11pt,twocolumn]{article}
\usepackage[margin=0.75in]{geometry}
\usepackage{graphicx}
\usepackage{siunitx}
\usepackage{amsmath}
\usepackage{float}
\usepackage{caption}
\usepackage{booktabs}
\usepackage{circuitikz}
\usepackage{url}


\title{Exploring Band Pass Behavior of Series RLC circuits}
\author{Jason Jain, Colby Hutchinson \\
        Oregon State University, Department of Physics \\
        Electronics Laboratory}
\date{March 3, 2026}

\graphicspath{{output/}}

\begin{document}

\twocolumn[
  \maketitle
  \begin{abstract}
    A series RLC circuit consists of an inductor $L$, a capacitor $C$, and a resistor $R$ in series, and has ubiquitous frequency filtering applications throughout electronics. We present a theoretical model for calculating amplitude and phase offset for an output signal from an RLC circuit, predicting strong attenuation within a band-pass region $\omega_0 \pm \beta$ of resonant angular frequency $\omega_0 = \qty{50}{\kilo\radian\per\second}$ and damping factor $\beta = \qty{1}{\kilo\radian\per\second}$. Driving the circuit with sinusoidal AC voltage across a range of input frequencies, we use an oscilloscope to measure experimental values for change in output amplitude and phase. We found the resonant frequency to be \qty{51.2}{\kilo\radian\per\second}, within 2.9\,\% of our theoretical prediction, and observed the expected attenuation behavior for phase and amplitude. Our theoretical model reasonably matches our measurements and successfully characterizes band-pass filtering behavior for series RLC circuits.
  \end{abstract}
  \vspace{1.5em}
]

\section{Introduction}\label{sec:introduction}
RLC circuits are fundamental components across everyday electronics, ubiquitous for their frequency filtering applications. Signals transmitting at the resonance frequency $\omega_0$ create a strong voltage response, which we characterize with admittance $Y = V_{\mathrm{out}}/(V_{\mathrm{in}} R)$. Frequencies at the band pass $\omega_0 \pm \beta$ fall by a factor of $|Y| = 1/\sqrt{2}$, while power drops by $P \propto V^2 \propto 1/2$; signals outside this range do not provide appreciable power and are thus filtered out. This band-pass filtering functionality has widespread applications for devices that communicate or operate at particular frequencies, like radios, WiFi chips~\cite{mo_wifi}, or power supplies~\cite{keysight}.

Admittance will be a primary quantity we analyze because it is independent of $V_{\mathrm{in}}$ (which may fluctuate). We'll also discuss phase shift $\phi_I$ between input/output current, important for calculating power delivery to electronic components. A phase shift of $\phi = 0$ means power is maximized, whereas a phase shift of $\phi = \pi/2$ means zero real power is delivered because the output is zero when the input is at a maximum.

We now examine an RLC circuit (Figure~\ref{fig:circuit}) with inductance $L = \qty{22}{\milli\henry}$, resistance $R = \qty{50}{\ohm}$, and capacitance $C = \qty{18}{\nano\farad}$. Sinusoidal AC voltage is our input, peaking at $|V_{\mathrm{in}}| \approx \qty{2}{\V}$, and we define output voltage across the resistor as $V_{\mathrm{out}} = V_R$.

\begin{figure}[H]
  \centering
  \begin{circuitikz}
    \draw (0, 0)
    to[V, v=$V_{E}$] (0, 3)
    to[R, l=\SI{50}{\ohm}, v=$v_R$] (3, 3)
    to[L, l=\SI{22}{\milli\henry}, v=$v_L$] (6, 3)
    to[C, l=\SI{18}{\nano\farad}, v=$v_C$] (6, 0)
    to[short, -*] (0, 0) node[ground]{};
  \end{circuitikz}
  \caption{The circuit analyzed in this lab consists of resistor $R$, inductor $L$, and capacitor $C$ in series with driving external voltage $V_E$. Voltage $V_E$ is a sine wave corresponding to input, and voltage $V_R$ corresponds to output.}
  \label{fig:circuit}
\end{figure}

\section{Theory}\label{sec:theory}



\begin{figure*}[t]
  \centering
  \includegraphics[width=\linewidth]{oscilloscope_traces.png}
  \caption{Traces of the real part of $V_E$ (Eq.~\eqref{eq:volt_phasor}) and $V_R$ (Eq.~\eqref{eq:vr_phasor}), assuming $|V_E| = 1$ for convenience. Starting from the top-left, we see a phase shift of $\phi_I \mp \pi/4$ within the band pass ($\omega \pm \beta$), where amplitudes drop to $1/\sqrt{2}$. In the top center at resonance ($\omega = \omega_0$), we see a phase shift of $\phi_I = 0$ and equal amplitudes---this is where the maximum response is achieved. In the bottom row, the phase shift approaches $\mp\pi/2$ as driving frequency $\omega$ approaches $\pm\infty$, while the amplitude approaches 0. We expect to see similar behavior in our experiment.}
  \label{fig:traces}
\end{figure*}


We model driving voltage $V_E$ as a complex phasor in Eq.~\eqref{eq:volt_phasor}, where only the real part corresponds to a physical quantity.
\begin{equation}
  \label{eq:volt_phasor}
  V_E = V_0 \, e^{i\omega t}
\end{equation}

We assume that power loss across the inductor and capacitor is negligible. Therefore, by Kirchhoff's Laws~\cite{gingrich} for Circuit~\ref{fig:circuit}, voltage across the inductor $V_L(t) = L \ddot{q}(t)$, resistor $V_R(t) = \dot{q}(t) R$, and capacitor $V_C(t) = q/C$ all sum to external voltage $V_E$.
\begin{equation}
  \label{eq:ode}
  \ddot{q} + \frac{R}{L} \dot{q} + \frac{1}{LC} q = \frac{V_0}{L} e^{i\omega t}
\end{equation}

Or reexpressed in standard driven harmonic oscillator form with damping factor $\beta = R/(2L) = \qty{1.1}{\kilo\radian\per\second}$ and resonant frequency $\omega_0 = (LC)^{-1/2} = \qty{50}{\kilo\radian\per\second}$,
\begin{equation}
  \label{eq:standard_form}
  \ddot{q} + 2\beta \dot{q} + \omega_0^2 q = \frac{V_0}{L} e^{i\omega t}
\end{equation}

\subsection{Transient Solutions}\label{sec:transient}
The characteristic equation of the homogeneous form of Eq.~\eqref{eq:standard_form} is $r^2 + 2\beta r + \omega_0^2 = 0$, with solutions given by the quadratic formula: $r_{1,2} = -\beta \pm \sqrt{\beta^2 - \omega_0^2}$. Since $\beta^2 < \omega_0^2$, we have $r_{1,2} = -\beta \pm i\sqrt{\omega_0^2 - \beta^2}$. This is an underdamped harmonic oscillator with a known transient solution~\cite{gingrich},
\begin{equation}
  \label{eq:transient}
  x(t) = A e^{-\beta t} \cos(\omega_1 t + \phi)
\end{equation}

The transient solution decays quickly enough for us to ignore it. Consider the damping time $t_D$ during which the circuit decays to $1/e$ of its starting value. Setting $\beta t_D = 1$ gives $t_D = 1/\beta = \qty{0.89}{\milli\s}$, during which we expect $\omega_1/(2\pi\beta) = 7$ cycles. This indicates that the transient solution decays almost immediately, so we can safely ignore it and focus on the steady-state solution.

\subsection{Steady State Solutions}\label{sec:steady}
The known steady-state solution~\cite{gingrich} to the ODE (Eq.~\eqref{eq:ode}) is
\begin{equation}
  \label{eq:amp_phase}
  \begin{aligned}
    q(t)     & = \frac{V_0/L \cdot e^{i\phi} e^{i\omega t}}{\sqrt{(\omega_0^2 - \omega^2)^2 + 4\beta^2\omega^2}}, \\
    \tan\phi & = \frac{-2\beta\omega}{\omega_0^2 - \omega^2}.
  \end{aligned}
\end{equation}

Differentiating $q(t)$ with respect to time gives current $I(t) = i\omega |q| e^{i\phi} e^{i\omega t} = \omega |q| e^{i(\phi+\pi/2)} e^{i\omega t}$. By Ohm's law, the voltage across the resistor is $V_R(t) = I(t) R$:
\begin{equation}
  \label{eq:vr_phasor}
  \begin{aligned}
    V_R(t) & = \frac{R\omega V_0/L}{\sqrt{(\omega_0^2 - \omega^2)^2 + 4\beta^2\omega^2}} e^{i(\omega t + \phi_I)}, \\
    \phi_I & = \phi + \frac{\pi}{2}.
  \end{aligned}
\end{equation}

We now derive theoretical admittance and phase from Eq.~\eqref{eq:vr_phasor}. Substituting external voltage $V_E = V_0 e^{i\omega t}$ and applying Ohm's Law, $V_R = I(t) R$ and $V_E = I(t) Z$ where $Z = 1/Y$ is the impedance, we solve for $Y$ and take the norm (since we can only measure magnitudes) to obtain the theoretical admittance (Eq.~\eqref{eq:y_th}).
\begin{equation}
  \label{eq:y_th}
  |Y|_{\mathrm{theoretical}}(\omega) = \frac{\omega/L}{\sqrt{(\omega_0^2 - \omega^2)^2 + 4\beta^2\omega^2}}
\end{equation}

Likewise, the theoretical phase shift $\phi_{I,\mathrm{th}}$ is given by Eqs.~\eqref{eq:amp_phase} and \eqref{eq:vr_phasor}:
\begin{equation}
  \label{eq:phi_th}
  \phi_{I,\mathrm{th}}(\omega) = \tan^{-1}\left[\frac{-2\beta\omega}{\omega_0^2 - \omega^2}\right] + \frac{\pi}{2}
\end{equation}

Within the band pass, we expect $\phi_I(\omega_0 \pm \beta) \approx \tan^{-1}\left[\frac{-2\beta\omega_0}{2\omega_0\beta}\right] + \frac{\pi}{2} = \pi/4$. Similarly, at the resonant frequency we expect $\phi_I(\omega_0) = 0$.

\section{Methods}\label{sec:methods}
We constructed Circuit~\ref{fig:circuit} on a breadboard with input sinusoidal voltage $V_E(t)$ from a function generator, peaking at $\approx \qty{2}{\V}$. Using an oscilloscope, we took measurements at 20 different driving frequencies between \qty{2}{\kilo\hertz} and \qty{25}{\kilo\hertz}, where $\omega = 2\pi f$.

Component and oscilloscope tolerances introduced minor uncertainty in our measurements. The voltage readings from our oscilloscope, particularly for $V_E$, flickered by $\pm 5\,\%$ during measurement. The nominal values of $R$, $L$, and $C$ are also subject to tolerances. Our measured results are given in Table~\ref{tab:data}.

To calculate experimental admittance (Eq.~\eqref{eq:y_exp}), we measured peak-to-peak amplitudes using the oscilloscope's built-in functionality. The factor of 2 from a peak-to-peak measurement cancels in the admittance ratio.
\begin{equation}
  \label{eq:y_exp}
  |Y|_{\mathrm{exp}}(\omega) = \frac{V_R}{V_E} \frac{1}{R}
\end{equation}

To calculate experimental phase (Eq.~\eqref{eq:phi_exp}), we measured the time shift $\Delta t$ between zero crossings of $V_E(t)$ and $V_R(t)$ using cursors on our oscilloscope.
\begin{equation}
  \label{eq:phi_exp}
  \phi_{I,\mathrm{exp}}(\omega) = 2\pi f \Delta t
\end{equation}

\section{Results}\label{sec:results}

\begin{table}[H]
  \centering
  \caption{Series RLC frequency response: external and resistor voltages, time shift, phase shift, and admittance. Frequency is scaled by $2\pi$ to convert to angular frequency.}
  \label{tab:data}
  \resizebox{\columnwidth}{!}{%
    \small
    \begin{tabular}{
        S[table-format=3.2]
        S[table-format=3.2]
        S[table-format=1.2]
        S[table-format=-2.2]
        S[table-format=-1.2]
        S[table-format=1.2]
      }
      \toprule
      {$\omega$ (\unit{\kilo\radian\per\second})} & {$V_R$ (\unit{mV})} & {$V_E$ (\unit{V})} & {$\Delta t$ (\unit{\micro\second})} & {$\phi_{I,\mathrm{exp}}$ (\unit{\radian})} & {$|Y|_{\mathrm{exp}}$ (\unit{\siemens})} \\
      \midrule
      12.57                                       & 28.00               & 2.12               & 124.00                              & 1.56                                       & 0.01                                     \\
      25.13                                       & 68.00               & 2.12               & 60.00                               & 1.51                                       & 0.03                                     \\
      31.42                                       & 98.00               & 2.12               & 48.00                               & 1.51                                       & 0.05                                     \\
      37.70                                       & 152.00              & 2.10               & 38.00                               & 1.43                                       & 0.07                                     \\
      43.98                                       & 284.00              & 2.04               & 30.00                               & 1.32                                       & 0.14                                     \\
      47.12                                       & 448.00              & 1.90               & 23.00                               & 1.08                                       & 0.24                                     \\
      49.01                                       & 576.00              & 1.68               & 17.00                               & 0.83                                       & 0.34                                     \\
      50.27                                       & 660.00              & 1.60               & 12.00                               & 0.60                                       & 0.42                                     \\
      50.89                                       & 680.00              & 1.50               & 6.00                                & 0.31                                       & 0.46                                     \\
      51.52                                       & 696.00              & 1.40               & 0.00                                & 0.00                                       & 0.48                                     \\
      52.15                                       & 688.00              & 1.50               & -3.00                               & -0.16                                      & 0.46                                     \\
      53.41                                       & 616.00              & 1.60               & -10.00                              & -0.53                                      & 0.38                                     \\
      54.66                                       & 544.00              & 1.80               & -15.00                              & -0.82                                      & 0.31                                     \\
      56.55                                       & 440.00              & 1.90               & -18.00                              & -1.02                                      & 0.23                                     \\
      59.69                                       & 320.00              & 2.00               & -20.00                              & -1.19                                      & 0.16                                     \\
      62.83                                       & 256.00              & 2.10               & -21.00                              & -1.32                                      & 0.12                                     \\
      69.12                                       & 172.00              & 2.10               & -21.00                              & -1.45                                      & 0.08                                     \\
      81.68                                       & 116.00              & 2.10               & -19.00                              & -1.55                                      & 0.05                                     \\
      113.10                                      & 60.00               & 2.10               & -13.20                              & -1.49                                      & 0.03                                     \\
      157.08                                      & 40.00               & 2.10               & -9.60                               & -1.51                                      & 0.02                                     \\
      \bottomrule
    \end{tabular}%
  }
\end{table}

We measured an experimental resonance frequency of $\omega = \qty{51.52}{\kilo\radian\per\second}$, where our phase shift is $\phi_I = \qty{0}{\radian}$ and our admittance is a maximum of $|Y|_{\mathrm{exp}} = \qty{0.48}{\siemens}$. This is very close to our theoretical resonance frequency of $\omega_0 = \qty{50}{\kilo\radian\per\second}$, with a percent difference of $\approx 2.9\,\%$.

As we move away from resonance, we see our phase shift approach $\pm\pi/2$, with data showing a phase shift of $\phi_I = \qty{-1.55}{\radian}$ at $\omega = \qty{81.68}{\kilo\radian\per\second}$ and $\phi_I = \qty{-1.49}{\radian}$ at $\omega = \qty{113.10}{\kilo\radian\per\second}$. We also see our admittance attenuate to zero, with voltage $V_R$ shrinking to $< \qty{50}{mV}$. These measurements align with our theoretical predictions in Figure~\ref{fig:traces}.

Our experimental admittance (Figure~\ref{fig:admittance}) is a reasonable fit to our theoretical curve. In addition to component tolerances, our fundamental assumptions introduce discrepancies to the model. Our calculation of Eq.~\eqref{eq:y_th} and Eq.~\eqref{eq:phi_th} assumed negligible power loss across the inductor and capacitor. When far from resonance, the circuit filters out most of the power and any power loss is negligible. Near resonance, however, the circuit approaches maximum power, so inefficiencies are more noticeable across the inductor and capacitor. It's unsurprising that the experimental amplitudes are much lower than theoretical predictions in this range.

Phase shift is a better fit than admittance for a broader frequency range, and can be used to calculate an experimental value for $\beta$. At $\omega = \omega_0 \pm \beta$, we should see a phase shift of $\phi_I = \pm\pi/4$. The closest data point we have is at $\omega = \qty{50.27}{\kilo\radian\per\second}$, where our phase shift is $\phi_I = \qty{0.83}{\radian}$, or within 5.4\,\% of $\pi/4$. Using the lower band-pass edge, we find an experimental $\beta$ of \qty{2.51}{\kilo\radian\per\second}, or within 54.7\,\% of our theoretical prediction of $\beta = \qty{1.1}{\kilo\radian\per\second}$. This is a more significant difference, but is likely the effect of component tolerances and model constraints, as we saw with admittance. Both the inductor and capacitor add small amounts of resistance, causing a higher effective $\beta = R_{\mathrm{eff}}/(2L)$ than predicted. These discrepancies are consistent with the expected effect of parasitic resistance in the inductor.

\begin{figure}[t]
  \centering
  \includegraphics[width=\columnwidth]{amplitude_plot.png}
  \caption{Admittance $|Y|$ vs.\ driving frequency $\omega$: experimental points and theoretical curve. Each sub-tick is $1\beta$ with the graph centered at $\omega_0$. We see a reasonable fit to our theoretical curve, with discrepancies expected from component tolerances and the fundamental assumptions of our model. The red line marks the resonant frequency $\omega_0$ where we see maximum response; the band-pass lines $\omega_0 \pm \beta$ intersect the curve at $|Y| = |Y|_{\max}/\sqrt{2}$, illustrating the narrow range of this circuit. Outside the band pass, power response drops by 1/2.}
  \label{fig:admittance}
\end{figure}

\begin{figure}[t]
  \centering
  \includegraphics[width=\columnwidth]{phase_plot.png}
  \caption{Phase $\phi_I$ vs.\ driving frequency $\omega$: experimental points and theoretical curve. Each sub-tick is $1\beta$ with the graph centered at $\omega_0$. We see a close fit to our theoretical curve, with discrepancies expected from component tolerances and model assumptions. The red line marks the resonant frequency $\omega_0$ where the phase shift is 0; the band-pass lines $\omega_0 \pm \beta$ intersect the curve at $\phi_I = \pm\pi/4$. Outside the band pass, the phase shift slowly approaches asymptotes of $\pm\pi/2$.}
  \label{fig:phase}
\end{figure}

\section{Conclusions}\label{sec:conclusions}

The RLC circuit is a basic oscillator circuit that is ubiquitous in electronics. When driven by a sinusoidal voltage, it behaves like a band pass filter with a resonant frequency $\omega_0$ and a bandwidth $\pm\beta$, attenuating signals outside this range. We built such an RLC circuit with $L = \qty{22}{\milli\henry}$, $R = \qty{50}{\ohm}$, and $C = \qty{18}{\nano\farad}$, and presented a theoretical model for filtered voltage amplitude and phase offset, dependent on driving frequency. We predicted a resonant frequency of $\omega_0 = \qty{50}{\kilo\radian\per\second}$ and sharp attenuation within the band-pass region $\omega_0 \pm \beta$ for $\beta = \qty{1.1}{\kilo\radian\per\second}$. Driving our circuit across a range of frequencies, we measured amplitude and phase shift data using an oscilloscope, and found an experimental $\omega_0 = \qty{51.2}{\kilo\radian\per\second}$, within 2.9\,\% of our theoretical prediction. Our theoretical curves fit reasonably well, with discrepancies expected from component tolerances and fundamental assumptions of our model. We assumed negligible power loss across the inductor and capacitor, an approximation which became less accurate near resonance, when power is maximized; therefore, we saw smaller-than-predicted amplitudes in this range. Observed attenuation behavior within the band-pass region matched our theoretical predictions, with phase shift reaching 0 at resonance and nearing $\pm\pi/4$ within the band-pass region. Our model gives engineers a theoretical framework for characterizing RLC circuits and predicting their behavior dependent solely on circuit parameters and driving frequency. By adjusting $L$ and $C$, one can tune $\omega_0$ to any desired frequency, while $R$ provides independent control over the bandwidth $\beta$, enabling precise frequency-selective filtering.

\begin{thebibliography}{9}
  \bibitem{gingrich}
  D.~M. Gingrich, \textit{Practical Electronics}, Chapter 2: Driven RLC Circuits. \url{https://www-f9.ijs.si/~gregor/ElektronikaVaje/ElectronicsLectures_GinGrich.pdf}

  \bibitem{copilot}
  GitHub Copilot, AI-powered code completion tool by GitHub. Used for assistance with Python code development and \LaTeX\ formatting. \url{https://github.com/features/copilot}

  \bibitem{keysight}
  Keysight Technologies. RLC circuit. \url{https://www.keysight.com/used/us/en/knowledge/glossary/oscilloscopes/rlc-circuit}

  \bibitem{mo_wifi}
  J.~Mo, ``Design and Characterization of a Three-layer PCB Low-pass Filter for Enhanced Harmonic Suppression in Wi-Fi Frequency Bands: An Application of Characteristic Mode Analysis,'' 2024 Photonics \& Electromagnetics Research Symposium (PIERS), Chengdu, China, 2024, pp.~1--4, doi: 10.1109/PIERS62282.2024.10618522.
\end{thebibliography}

\section*{Addendum}
The native language of the author(s) is English.

\end{document}
